% ===================================================================
%  TABULKA č. 13 — Hotel. — Restaurace. Kavárna.
%  Traduction 2026 d'apres texto/io, controlee sur texto/fr, texto/ru
%  et texto/uk. Consigne : voir texto/cs/10-tabulka-01.tex.
%
%  DEUX NOTES D'AUTEUR, l'une et l'autre en gras-italique dans le
%  fac-simile, comme dans l'ido : la premiere renvoie au tableau 6
%  pour le detail de la chambre, la seconde au vocabulaire pour la
%  carte des plats.
%
%  LE TABLEAU DE L'HOTEL PARTAGE SON LEXIQUE EN DEUX A L'INTERIEUR
%  D'UNE MEME PIECE, quatrieme cas apres la montagne, la mer et
%  l'hotel galicien. Ce qui SE MANGE est tcheque de bout en bout —
%  « máslo » le beurre, « sýr » le fromage, « ryba » le poisson,
%  « ovoce » les fruits, « špenát », « kýta » le gigot — et ce qui SE
%  SERT arrive avec le service : « hotel », « restaurace »,
%  « výtah », « telegram », « biliár », « automobil », « omnibus ».
%
%  ET « SPROPITNÉ » N'EST NI « PROPINA » NI « POURBOIRE ». Le
%  galicien, l'espagnol et le portugais gardaient ensemble le latin
%  de ce qu'on boit en plus ; le tcheque a calque l'allemand
%  « Trinkgeld » — l'argent a boire — mot pour mot, avec ses propres
%  racines. C'est le troisieme etat de la meme idee : garder le mot,
%  le refaire, ou le traduire piece a piece.
% ===================================================================

%%K t13-noto-1 noto
\VUnotes{143.2mm}{%
\textit{\VUgras{(*) O podrobnostech pokoje viz tabulku č. 6.}}%
}

%%K t13-apar-1 apar
\VUsaut{3.0mm}
\VUtitre{15.0pt}{60}{TABULKA č. 13}
\VUsaut{1.6mm}
\VUpk{10.2pt}{40}{Hotel. --- Restaurace.}
\VUsaut{2.0mm}
\VUpk{10.2pt}{40}{Kavárna.}
\VUsaut{3.4mm}

%%K t13-01-1 p
1. --- Když jsem včera večer přijel do Royanu, ubytoval jsem se v «\textit{Hotelu Oceán}»,
který mi doporučil jeden přítel.

%%K t13-01-2 p
Dali mi krásný \VUgras{pokoj} \textsuperscript{(2)} ve druhém \VUgras{patře}
\textsuperscript{(1)}, kde jsem se velmi dobře vyspal (*).

%%K t13-02-1 p
2. --- Dnes ráno, když jsem vycházel ze svého pokoje, kam mi \VUgras{služebná}
\textsuperscript{(3)} přinesla snídani, vstoupil jsem do \VUgras{kuřárny}
\textsuperscript{(5)}, které se užívá také jako \VUgras{čítárny} \textsuperscript{(4)}, abych
si přečetl \VUgras{noviny} \textsuperscript{(6)} a zakouřil si \VUgras{cigaretu}
\textsuperscript{(8)}. Po čtení jsem sjel
\VUgras{výtahem} \textsuperscript{(9)} a šel se projít městem.

%%K t13-03-1 p
3. --- Protože jsem zapomněl zeptat se hotelového \VUgras{účetního} \textsuperscript{(12)}, v
kolik hodin se podávají jídla u \VUgras{společného stolu} \textsuperscript{(11)}, vrátil jsem
se brzy a využil toho, abych se šel vykoupat do hotelové \VUgras{koupelny}
\textsuperscript{(10)}. Potom jsem šel znovu k
\VUgras{pokladně} \textsuperscript{(15)}, abych zatelefonoval jednomu ze svých přátel.
Ale \VUgras{telefon} \textsuperscript{(13)} byl obsazen; když \VUgras{pokladní}
\textsuperscript{(14)}, která zapisovala \VUgras{účet} \textsuperscript{(17)} do
\VUgras{knihy hostů} \textsuperscript{(16)}, viděla mé rozpaky, požádala
\VUgras{vrátného} \textsuperscript{(18)}, aby poslal \VUgras{poslíčka}
\textsuperscript{(19)} vyřídit můj vzkaz \VUgras{na kole} \textsuperscript{(20)}.

%%K t13-04-1 p
4. --- Poděkoval jsem jí a vyšel ven, abych dal spropitné \VUgras{nosiči}
\textsuperscript{(24)} (nádeníkovi), který přivezl z nádraží na svém \VUgras{ručním vozíku}
\textsuperscript{(25)} mou \VUgras{krabici na klobouky} \textsuperscript{(26)},
\VUgras{můj kufřík} \textsuperscript{(27)} a \VUgras{mou cestovní tašku}
\textsuperscript{(28)}. \VUgras{Listonoš} \textsuperscript{(21)}, který právě přišel, mi dal
\VUgras{dopis} \textsuperscript{(22)}.

%%K t13-05-1 p
5. --- Zůstal jsem ještě chvíli na prahu dveří, blízko \VUgras{poštovní schránky}
\textsuperscript{(23)}, a díval se na ruch na ulici. \VUgras{Kočí} \textsuperscript{(30)}
\VUgras{omnibusu} \textsuperscript{(29)} nakládal na svůj vůz
\VUgras{kufr} \textsuperscript{(31)} \VUgras{cestujícího} \textsuperscript{(32)},
kterého doprovázel \VUgras{hoteliér} \textsuperscript{(33)}. Mladý
\VUgras{telegrafista} \textsuperscript{(35)} bez spěchu přinášel \VUgras{telegram}
\textsuperscript{(34)} pro hotelového hosta; \VUgras{turista} \textsuperscript{(36)} s
\VUgras{fotografickým přístrojem} \textsuperscript{(38)} přes rameno nahlížel do svého
\VUgras{průvodce} \textsuperscript{(37)}.

%%K t13-tit-1 sub
\VUsaut{4.0mm}
\VUpk{10.2pt}{40}{Čas oběda.}
\VUsaut{3.0mm}

%%K t13-06-1 p
6. --- Před mříží bezstarostný \VUgras{poslíček} \textsuperscript{(39)} podřimoval mezi svým
\VUgras{nosným hákem} \textsuperscript{(40)} a svou \VUgras{krabicí s krémem na boty}
\textsuperscript{(41)}, zatímco mladá \VUgras{pradlena} \textsuperscript{(42)}, která nesla
\VUgras{koš} \textsuperscript{(43)} prádla, se smála vážné tváři
\VUgras{vrchního} \textsuperscript{(44)}, který stál u dveří.

%%K t13-07-1 p
7. --- Když jsem uviděl \VUgras{kuchaře} \textsuperscript{(45)}, který vyšel ze své
\VUgras{kuchyně} \textsuperscript{(47)} položené v \VUgras{suterénu}
\textsuperscript{(48)}, aby poslal \VUgras{kuchtíka} \textsuperscript{(46)} vyřídit vzkaz,
vzpomněl jsem si, že hodina oběda je blízko, a vstoupil jsem do restaurace, procházeje dvorem,
v němž \VUgras{řidič} \textsuperscript{(50)} odstavoval svůj
\VUgras{automobil} \textsuperscript{(49)}, zatímco \VUgras{mechanik}
\textsuperscript{(52)} opravoval, podle pokynů \VUgras{cyklisty} \textsuperscript{(53)},
\VUgras{benzinovou tříkolku} \textsuperscript{(51)}, která právě měla nehodu.

%%K t13-08-1 p
8. --- Zeptal jsem se mladého \VUgras{poslíčka} \textsuperscript{(54)}, který nesl
\VUgras{sklenice piva} \textsuperscript{(55)}, zda je společný stůl pod verandou, a na
jeho kladnou odpověď jsem zaujal místo u jednoho ze stolů a dal jsem si podat výborný oběd.

%%K t13-noto-2 noto
\VUnotes{143.2mm}{%
\textit{\VUgras{(*) S pomocí seznamu jídel zavedeného ve slovníku bude moci učitel snadno obměňovat níže uvedený jídelní lístek.}}%
}

%%K t13-tit-2 sub
\VUsaut{4.0mm}
\VUpk{10.2pt}{40}{Výborný oběd.}
\VUsaut{3.0mm}

%%K t13-09-1 p
9. --- Číšník mi přinesl \VUgras{jídelní lístek} (*) oběda a nápojový lístek. Vybral jsem si a
objednal jako \VUgras{předkrm krevety} s \VUgras{máslem,} a jako
\VUgras{první chod dršťky po caenském způsobu.} Nejedl jsem \VUgras{rybu,} ale požádal
jsem o porci skopové \VUgras{kýty} a jako \VUgras{zeleninový chod špenát} se smetanou. Potom,
protože se mi žádný z \VUgras{mezichodů} nelíbil, dal jsem si přinést různé zákusky,
\VUgras{sýr,} \VUgras{ovoce} a \VUgras{sušenky.} Přidal jsem ještě výborné
\VUgras{víno} Saint-Émilion a, velmi spokojen se svým jídlem, opustil jsem stůl, když
jsem dal \VUgras{číšníkovi} \textsuperscript{(57)} pěkné \VUgras{spropitné}
\textsuperscript{(59)}.

%%K t13-10-1 p
10. --- Když viděl, že jsem připraven odejít, navrhl mi \VUgras{správce}
\textsuperscript{(60)}, abych šel na balkon obdivovat nádherné panorama
\VUgras{oceánu} \textsuperscript{(62)}. V dálce se objevovaly rybářské
\VUgras{bárky} \textsuperscript{(63)}, \VUgras{plachetnice} \textsuperscript{(64)} a
obrovský \VUgras{parník} \textsuperscript{(65)}, jehož černá silueta vystupovala na bílých
\VUgras{oblacích} \textsuperscript{(67)} \VUgras{obzoru} \textsuperscript{(66)}. Lehký vánek
rozvlnil \VUgras{prapor} \textsuperscript{(70)} zasazený nad \VUgras{znakem}
\textsuperscript{(69)} \VUgras{haly} \textsuperscript{(68)}.

%%K t13-tit-3 sub
\VUsaut{3.0mm}
\VUpk{10.2pt}{40}{Zábavy.}
\VUsaut{3.0mm}

%%K t13-11-1 p
11. --- Podívaná byla tak zajímavá, že jsem se rozhodl vystoupit zítra na terasu kavárny a
vzít s sebou \VUgras{divadelní kukátko} \textsuperscript{(71)} nebo
\VUgras{dalekohled} \textsuperscript{(72)}.

%%K t13-12-1 p
12. --- Když jsem opustil balkon, šel jsem do hotelové kavárny vypít \VUgras{nápoj}
\textsuperscript{(73)} a zahrát si \VUgras{partii biliáru} \textsuperscript{(75)}. Bez
zastávky jsem prošel kavárenským sálem, v němž mnoho hostů hrálo \VUgras{šachy}
\textsuperscript{(78)}, \VUgras{dámu} \textsuperscript{(79)}, \VUgras{domino}
\textsuperscript{(80)}, \VUgras{vrhcáby} \textsuperscript{(81)} nebo \VUgras{karty}
\textsuperscript{(82)}, a šel jsem přímo nahoru do \VUgras{biliárového sálu}
\textsuperscript{(74)}.

%%K t13-13-1 p
13. --- Když partie skončila, sešel jsem do přízemí a požádal jsem číšníka o
\VUgras{obálku} \textsuperscript{(84)}, \VUgras{papír} \textsuperscript{(83)},
\VUgras{poštovní známku} \textsuperscript{(85)}, \VUgras{pero}
\textsuperscript{(87)} a \VUgras{inkoust} \textsuperscript{(86)}, abych napsal dopis.

%%K t13-13-2 p
Potom jsem zavolal \VUgras{drožkáře} \textsuperscript{(92)}, jehož \VUgras{vůz}
\textsuperscript{(88)} zastavil před kavárnou. Zeptal jsem se ho na cenu podle hodin nebo za
jízdu, a když jsme se shodli na ceně, otevřel mi \VUgras{dvířka vozu} \textsuperscript{(89)} a
usadil mě. Vystoupil znovu na \VUgras{kozlík} \textsuperscript{(91)}, rychle pobídl koně
\VUgras{bičíkem} \textsuperscript{(93)} a vyjeli jsme na půvabnou projížďku.
